%==============================================================================
%  Standalone note: Pressure-mean drift under the iterated (Codina) penalty
%  in the all-Dirichlet setting, and a re-centering PROPOSAL.
%
%  Notation is kept compatible with the manuscript
%  "A stabilized finite element method for incompressible, inertial flows
%   in inhomogeneous porous media" (Casas, Gonzalez-Usua, Codina, de-Pouplana).
%
%  STATUS: exploratory. The re-centering step described in Section 6 is a
%  PROPOSAL of this note, NOT part of Codina's published iterated-penalty
%  method and NOT (yet) implemented in the solver. See Section 7 (Provenance).
%==============================================================================
\documentclass[11pt]{article}

\usepackage[utf8]{inputenc}
\usepackage[margin=1in]{geometry}
\usepackage{amsmath,amssymb,amsthm,mathtools}
\usepackage{bm}
\usepackage{xcolor}
\usepackage[colorlinks=true,allcolors=blue]{hyperref}
\usepackage{cleveref}

% Theorem environments
\theoremstyle{plain}
\newtheorem{proposition}{Proposition}
\newtheorem{lemma}{Lemma}
\theoremstyle{remark}
\newtheorem{remark}{Remark}

% Shorthands local to this note
\newcommand{\bu}{\boldsymbol{u}}
\newcommand{\bv}{\boldsymbol{v}}
\newcommand{\bn}{\boldsymbol{n}}
\newcommand{\bx}{\boldsymbol{x}}
\newcommand{\uh}{\boldsymbol{u}_h}
\newcommand{\ph}{p_h}
\newcommand{\qh}{q_h}
\newcommand{\Om}{\Omega}
\newcommand{\dOm}{\partial\Omega}
\newcommand{\meas}{\lvert\Omega\rvert}
\newcommand{\enum}{\varepsilon_{\mathrm{num}}}
\newcommand{\ephys}{\varepsilon_{\mathrm{phys}}}
\newcommand{\pbar}{\overline{p}}
\newcommand{\norm}[1]{\left\Vert #1\right\Vert}
\newcommand{\abs}[1]{\left\lvert #1\right\rvert}
\newcommand{\dx}{\,\mathrm{d}\Omega}
\newcommand{\ds}{\,\mathrm{d}s}
\newcommand{\Ltwo}{L^2(\Omega)}

\title{Pressure-mean drift under the iterated penalty,\\[2pt]
and a re-centering proposal (exploratory)}
\author{}
\date{}

\begin{document}
\maketitle

\begin{abstract}
The all-Dirichlet velocity boundary condition leaves the discrete pressure determined only up to an
additive constant. In the solver this gauge is fixed \emph{weakly} by the Codina iterated
(``numerical'') penalty $\enum(p^{n}-p^{n-1})$ in the continuity equation, which pins the constant to
the \emph{previous} outer iterate rather than to zero mean. We show that, because the discrete boundary
flux $\rho=\oint_{\dOm}\alpha\,\uh\cdot\bn\ds$ is a fixed nonzero residue at the interpolation/quadrature
level, each outer penalty pass shifts the pressure mean by a fixed amount
$-\rho/(\enum\meas)$, so the mean drifts linearly in the pass count while the usual relative-drift
stopping test converges \emph{spuriously} through a growing denominator. We then describe a
\textbf{re-centering proposal}: subtract the domain mean of $p$ between passes. Because the constant is
pure gauge, this is exact and changes no mean-removed quantity, while it bounds the mean and makes the
stopping test honest. We are careful to separate what is established from what is proposed
(\cref{sec:provenance}): the iterated penalty is Codina's / the manuscript's; the \emph{re-centering}
is a proposal of this note, which we have \emph{not} found in Codina's work. It is implemented behind a
config flag (\texttt{recenter\_pressure\_between\_penalty\_passes}, default \emph{off}) and A/B-verified
(\cref{sec:empirical}): with the pressure constant free, every mean-removed error is bit-identical
on/off, while the stored pressure mean collapses from $\mathcal{O}(10^{2})$ to machine zero. Note that
what is dormant is the multi-pass \emph{accumulation}, not the mean itself: on a coarse $P_1$ mesh the
\emph{single}-pass shift already reaches $\approx 230$ and dominates $\norm{p}$, so re-centering the
returned field is materially useful, not merely defensive.
\end{abstract}

\section{Setting}
\label{sec:setting}

On a bounded polyhedral domain $\Om\subset\mathbb{R}^{d}$ we consider the stabilized porous
Navier--Stokes system of the manuscript, with velocity $\uh\in V_h$ and pressure $\ph\in Q_h$ on an
equal-order Lagrangian pair. The physics is \emph{incompressible}: there is no genuine compressibility,
so the physical penalty $\ephys=0$. The velocity is prescribed on the \emph{entire} boundary,
$\uh=\bu_{\mathrm{ex}}$ on $\dOm$ (all-Dirichlet, as in the manuscript's \S5.2 three-dimensional MMS).

\begin{lemma}[Pressure gauge freedom]
\label{lem:gauge}
Under all-Dirichlet velocity data the pressure enters the momentum balance only through the pairing
$-\bigl(p,\ \alpha\,\nabla\!\cdot\bv+\bv\cdot\nabla\alpha\bigr)=-\bigl(p,\ \nabla\!\cdot(\alpha\bv)\bigr)$,
which for any admissible test field $\bv$ vanishing on $\dOm$ annihilates constants:
$\int_{\Om}\nabla\!\cdot(\alpha\bv)\dx=\oint_{\dOm}\alpha\,\bv\cdot\bn\ds=0$. Hence a constant added to
$p$ leaves the momentum residual unchanged, and the continuous solution is unique only up to that
constant.
\end{lemma}

The variational continuity equation used by the solver, tested against $q\in Q_h$, is
\begin{equation}
\label{eq:mass}
\bigl(q,\ \ephys\,p^{n}+\nabla\!\cdot(\alpha\,\bu^{n})\bigr)
\;+\;\underbrace{\enum\,\bigl(q,\ p^{n}-p^{n-1}\bigr)}_{\text{iterated penalty}}
\;+\;S_q(q;\bu^{n},p^{n})
\;=\;\bigl(q,\ g\bigr),
\end{equation}
where $g$ is the (manufactured) mass source, $S_q$ collects the VMS subscale contributions of the
pressure-test block ($\tau_2$/orthogonal-projection couplings), and $\nabla\!\cdot(\alpha\bu)
=\alpha\,\nabla\!\cdot\bu+\bu\cdot\nabla\alpha$. The superscript $n$ is the outer penalty-pass index:
$p^{n-1}$ is held \emph{fixed} within pass $n$, and updated $p^{n-1}\leftarrow p^{n}$ between passes.

\section{What the iterated penalty pins}
\label{sec:penalty}

Equation~\eqref{eq:mass} carries $\enum p^{n}$ on the left and $\enum p^{n-1}$ lagged to the right, so the
net residual contribution is $\enum(p^{n}-p^{n-1})$. This is the classical \emph{iterated-penalty}
(equivalently augmented-Lagrangian / Uzawa) device for the incompressibility constraint: it regularizes
the otherwise singular pressure block \emph{without} driving a penalty parameter to zero, and it
vanishes at convergence ($p^{n}=p^{n-1}$), so the incompressible solution is recovered unaltered. In the
present, gauge-free setting its second role is to render the discrete system nonsingular by pinning the
constant-pressure null mode of \cref{lem:gauge}.

\begin{remark}[The pin is to $p^{n-1}$, not to zero mean]
\label{rem:pin}
Crucially, $\enum(p^{n}-p^{n-1})$ pins $p^{n}$ toward the constant carried by $p^{n-1}$. It does
\emph{not} impose $\int_{\Om}p^{n}\dx=0$. Zero-mean is applied only in post-processing, when the error
norm $\norm{p-p_{\mathrm{ex}}}_{\Ltwo}$ is evaluated with the mean removed. The solved/stored pressure
field itself is never re-centered.
\end{remark}

\section{The mean-drift mechanism}
\label{sec:drift}

\begin{proposition}[Per-pass mean shift]
\label{prop:drift}
Let $\ephys=0$ and let the constant mode $q\equiv1$ be admissible in $Q_h$. Assume the pressure-test
subscale form is orthogonal to constants, $S_q(1;\cdot,\cdot)=0$ \textup{(}exact for OSGS, where the
projected residual is $L^2$-orthogonal to $Q_h\ni\mathrm{const}$; higher order for
ASGS\textup{)}. Then the pass-$n$ pressure mean $\pbar^{\,n}:=\meas^{-1}\int_{\Om}p^{n}\dx$ satisfies
\begin{equation}
\label{eq:shift}
\pbar^{\,n}-\pbar^{\,n-1} \;=\; -\,\frac{\rho^{n}-\int_{\Om} g\dx}{\enum\,\meas},
\qquad
\rho^{n}:=\oint_{\dOm}\alpha\,\bu^{n}\cdot\bn\ds .
\end{equation}
In particular, for a solenoidal manufactured field \textup{(}$\nabla\!\cdot(\alpha\,\bu_{\mathrm{ex}})=0$,
hence $\int_\Om g\dx=0$\textup{)} with a converged velocity, $\rho^{n}\to\rho$ is a fixed constant and the
mean advances by the same increment $-\rho/(\enum\meas)$ \emph{every} pass:
$\pbar^{\,n}=\pbar^{\,0}-n\,\rho/(\enum\meas)$.
\end{proposition}

\begin{proof}
Test \eqref{eq:mass} with $q\equiv1$. The penalty term gives
$\enum\int_{\Om}(p^{n}-p^{n-1})\dx=\enum\meas(\pbar^{\,n}-\pbar^{\,n-1})$. With $\ephys=0$ and
$S_q(1;\cdot,\cdot)=0$, the Galerkin term is
$\int_{\Om}\nabla\!\cdot(\alpha\bu^{n})\dx=\oint_{\dOm}\alpha\bu^{n}\cdot\bn\ds=\rho^{n}$ by the
divergence theorem, and the right-hand side is $\int_\Om g\dx$. Solving for the mean increment gives
\eqref{eq:shift}.
\end{proof}

The residue $\rho$ is nonzero only because $\bu^{n}$ is the discrete (interpolant-level) velocity: the
exact field satisfies $\oint_{\dOm}\alpha\,\bu_{\mathrm{ex}}\cdot\bn\ds=\int_\Om\nabla\!\cdot(\alpha
\bu_{\mathrm{ex}})\dx=0$ identically, so $\rho$ lives at the interpolation/quadrature error level. It is
small, but it is \emph{fixed and generically nonzero}, which is exactly what produces an unbounded,
h-\emph{independent} accumulation of the mean.

\section{Spurious convergence of the relative-drift stopping test}
\label{sec:stop}

The outer loop stops on the relative pressure drift
\begin{equation}
\label{eq:driftstop}
\delta^{n} \;=\; \frac{\norm{p^{n}-p^{n-1}}_{\Ltwo}}{\norm{p^{n}}_{\Ltwo}} \;<\; \mathrm{xtol}.
\end{equation}
Once the pressure \emph{shape} has stabilized, the numerator is dominated by the constant shift,
$\norm{p^{n}-p^{n-1}}_{\Ltwo}\approx\meas^{1/2}\abs{\rho}/(\enum\meas)$, a constant, while the
denominator grows with the running mean, $\norm{p^{n}}_{\Ltwo}\gtrsim\meas^{1/2}\abs{\pbar^{\,n}}
\sim\meas^{1/2}\,n\,\abs{\rho}/(\enum\meas)$. Hence
\begin{equation}
\delta^{n}\;\sim\;\frac{1}{n}\;\xrightarrow[n\to\infty]{}\;0 ,
\end{equation}
so \eqref{eq:driftstop} is eventually satisfied \emph{not} because the pressure has converged but
because the denominator diverges. Tightening $\mathrm{xtol}$ makes this \emph{worse}: it forces more
passes and hence a larger accumulated mean.

\section{Consequences}
\label{sec:consequences}

\begin{enumerate}
\item[(i)] \textbf{Loss of relative precision in $p$.} Once $\abs{\pbar}\gg\abs{p-\pbar}$, the physically
meaningful fluctuation is stored in the low-order digits of a large floating-point number; the
mean-removed error metric is immune, but any downstream use of the raw field degrades.
\item[(ii)] \textbf{Quadrature-level momentum contamination.} By \cref{lem:gauge} a constant $p_0$
cancels from momentum \emph{analytically}, $\int_{\Om}p_0\,\nabla\!\cdot(\alpha\bv)\dx=0$. Discretely
this cancellation is only as exact as the quadrature of $\nabla\!\cdot(\alpha\bv_h)$ for variable
$\alpha$, so a large $p_0$ injects a spurious body force $\mathcal{O}\!\bigl(p_0\times\text{(quadrature
error)}\bigr)$ into the velocity equation.
\item[(iii)] \textbf{A dishonest stopping test} (\cref{sec:stop}).
\end{enumerate}

\section{The re-centering proposal}
\label{sec:fix}

\begin{quote}
\emph{Proposal.} \textbf{(a)} Between outer passes, subtract the domain mean before storing the lagged
pressure, $\;p^{n-1}\leftarrow p^{n}-\pbar^{\,n}$; \textbf{(b)} re-center the \emph{returned} pressure
once at the end, $p\leftarrow p-\pbar$.
\end{quote}

Both parts matter, and for different reasons. Part (a) prevents the multi-pass \emph{accumulation}
(\cref{prop:bounded}); but by \cref{prop:drift} even a single pass from a zero-mean lag lands at
$\pbar\approx-\rho/(\enum\meas)$, which on a coarse mesh already dominates $\norm{p}$ (the A/B below finds
$\pbar\approx230$ against a true $\norm{p}\approx0.7$). Part (b) therefore does the work of curing
consequences (i)--(ii): it drives the \emph{stored/exported} field to zero mean. By \cref{lem:gauge} the
additive constant is pure gauge, so both parts leave the velocity, the pressure \emph{shape}, and every
mean-removed quantity exactly unchanged:

\begin{proposition}[Bounded mean and an honest test]
\label{prop:bounded}
With re-centering, $\pbar^{\,n-1}=0$ entering each pass, so by \eqref{eq:shift} the mean never
accumulates: $\abs{\pbar^{\,n}}\le\abs{\rho}/(\enum\meas)$ for all $n$. Moreover the drift numerator in
\eqref{eq:driftstop} then measures the genuine change of the pressure \emph{shape} (the constant having
been removed from both operands), so $\delta^{n}\to0$ reflects convergence rather than denominator
growth.
\end{proposition}

\Cref{prop:bounded} also upgrades the stopping test from a heuristic that passes for the wrong reason to
one that certifies the quantity of interest.

\section{Provenance: what is Codina's and what is proposed here}
\label{sec:provenance}

This section is deliberately explicit, because the two ingredients have very different pedigrees.

\begin{itemize}
\item \textbf{The iterated penalty $\enum(p^{n}-p^{n-1})$ is established.} It is the device used in the
manuscript (\S5.2) and belongs to the classical iterated-penalty / augmented-Lagrangian family for the
incompressibility constraint (e.g.\ Fortin \& Glowinski, \emph{Augmented Lagrangian Methods}, 1983; and
the stabilized-FEM line associated with Codina). Its convergence-as-$p^{n}\!\to\!p^{n-1}$ and its role in
regularizing the pressure block are standard.

\item \textbf{The re-centering of \cref{sec:fix} is a proposal of this note, and we do \emph{not}
attribute it to Codina.} We have not located a between-pass mean-subtraction step in Codina's
iterated-penalty presentations, and it should be read as an original hardening pending a literature
check. Two points make its absence unsurprising rather than suspicious:
\begin{enumerate}
\item[(a)] The mean-drift pathology is specific to the \emph{pure-Dirichlet-velocity} gauge-free setting
(\cref{lem:gauge}). Applications that carry a pressure/traction boundary condition, or that impose the
zero-mean constraint directly (a Lagrange multiplier, a fixed pressure DOF, or a null-space projection),
never expose the drift, so the question of re-centering the \emph{lag} does not arise there.
\item[(b)] Where the constant is fixed by other means, the iterated penalty is applied to a pressure that
is \emph{already} gauge-fixed each pass; re-centering the lag is then the natural transcription of that
practice to a solver whose \emph{only} gauge control is the penalty itself. In other words, the proposal
is arguably what one \emph{implicitly} assumes when the constant is otherwise pinned --- but that
assumption is not part of the bare iterated-penalty recipe, and asserting it as ``Codina's'' would be
incorrect.
\end{enumerate}
\end{itemize}

\noindent
The honest status is therefore: \emph{iterated penalty --- established and in use; between-pass
re-centering --- our proposal, unverified against the source literature, not implemented.}

\section{Empirical status: accumulation dormant, but the mean is not small}
\label{sec:empirical}

Two facts must be kept apart. First, the multi-pass \emph{accumulation} is dormant: the outer loop reaches
a genuine fixed point in two passes,
\[
\text{pass 1: } \delta^{1}\approx0.8\text{--}1.0,\qquad \text{pass 2: } \delta^{2}=0.0,
\]
i.e.\ $p^{2}=p^{1}$ (the penalty term vanishes), so the mean shifts at most once and never runs away.
Second, and less benign, that \emph{single} shift is not small: $\rho$ is a genuine interpolation/
quadrature residue that grows on coarser meshes and lower order, and by \eqref{eq:shift} the resulting
mean is $\rho/(\enum\meas)$ with the small $\enum\approx1.7\cdot10^{-5}$ in the denominator.

An A/B on the gauge-free ($\ephys=0$) case --- ASGS $P_1$ on the structured Kuhn $(8,8,2)$ mesh, flag off
vs.\ on --- makes both points concrete:
\begin{center}
\begin{tabular}{lccccc}
\hline
re-center & $L^2u$ (mean-rm) & $L^2p$ (mean-rm) & $H^1u$ & $\overline{p}$ (raw) & $\norm{p}$ (raw)\\
\hline
OFF & $2.5293\!\cdot\!10^{-2}$ & $3.5022\!\cdot\!10^{-1}$ & $3.1272\!\cdot\!10^{-1}$ & $2.30\!\cdot\!10^{2}$ & $1.26\!\cdot\!10^{2}$\\
ON  & $2.5293\!\cdot\!10^{-2}$ & $3.5022\!\cdot\!10^{-1}$ & $3.1272\!\cdot\!10^{-1}$ & $-1.2\!\cdot\!10^{-13}$ & $6.98\!\cdot\!10^{-1}$\\
\hline
\end{tabular}
\end{center}
The mean-removed errors are \emph{bit-identical} (behavior-preserving), confirming the gauge argument. But
raw $\overline{p}$ collapses from $230$ to machine zero, and $\norm{p}$ from $126$ to $0.70$: with the flag
off, the physical pressure fluctuation ($\approx0.7$) sits buried under a mean of $230$ --- exactly
consequence~(i). Re-centering is therefore a genuine improvement to the stored/exported field, not merely
defensive; it is most valuable on coarse/low-order meshes, non-solenoidal target flows, or any use of the
raw pressure, and it makes the drift stopping test honest everywhere.

\section{Drawbacks and where NOT to apply it}
\label{sec:drawbacks}

The proposal has no drawback \emph{within its hypotheses}, but those hypotheses matter:

\begin{itemize}
\item \textbf{Re-center only when the constant is genuinely free.} Under \cref{lem:gauge} (all-Dirichlet
velocity, $\ephys=0$) the mean is arbitrary and removing it is exact. If instead the problem carries a
Neumann/traction boundary, an open outlet, or a genuine net mass exchange, the pressure level is
\emph{physical}; subtracting its mean each pass would then silently annihilate a real quantity and could
mask a true mass imbalance. Any implementation must therefore be gated on the pure-Dirichlet,
gauge-free case, not applied unconditionally.
\item \textbf{It does not touch accuracy.} By construction it changes no mean-removed error; it must not
be oversold as improving convergence rates or fixing any discretization defect (in particular it is
unrelated to the P2 element-family coercivity question).
\item \textbf{It is not the only cure.} A hard zero-mean constraint (Lagrange multiplier / pinned DOF /
null-space deflation) achieves the same gauge fixing more directly; the iterated penalty exists precisely
to \emph{avoid} such a hard constraint, and re-centering is the lightweight transcription consistent with
that choice. Capping the number of passes is a cruder alternative that bounds but does not eliminate the
accumulation.
\end{itemize}

\section*{Implementation status}
\emph{Implemented behind a default-off flag.} Both parts of \cref{sec:fix} are wired into the outer
penalty loop under \texttt{recenter\_pressure\_between\_penalty\_passes} (default \texttt{false}
$\Rightarrow$ bit-identical legacy). \texttt{validate!} enforces the two applicability guards of
\cref{sec:drawbacks}: it requires \texttt{iterative\_penalty\_enabled} and $\ephys=0$ (\texttt{eps\_val}
$=0$). The A/B of \cref{sec:empirical} confirms mean-removed errors are bit-identical on/off while the
stored mean is driven to machine zero. What remains \emph{pending} is the editorial decision of whether to
adopt it by default and how to cite it --- gated by the provenance question of \cref{sec:provenance} (it
is our proposal, not Codina's).

\end{document}
